\IfFileExists{stacks-project.cls}{%
\documentclass{stacks-project}
}{%
\documentclass{amsart}
}

% For dealing with references we use the comment environment
\usepackage{verbatim}
\newenvironment{reference}{\comment}{\endcomment}
%\newenvironment{reference}{}{}
\newenvironment{slogan}{\comment}{\endcomment}
\newenvironment{history}{\comment}{\endcomment}

% For commutative diagrams we use Xy-pic
\usepackage[all]{xy}

% We use 2cell for 2-commutative diagrams.
\xyoption{2cell}
\UseAllTwocells

% We use multicol for the list of chapters between chapters
\usepackage{multicol}

% This is generally recommended for better output
\usepackage{lmodern}
\usepackage[T1]{fontenc}

% For cross-file-references
\usepackage{xr-hyper}

% Package for hypertext links:
\usepackage{hyperref}

% For any local file, say "hello.tex" you want to link to please
% use \externaldocument[hello-]{hello}
\externaldocument[introduction-]{introduction}
\externaldocument[conventions-]{conventions}
\externaldocument[sets-]{sets}
\externaldocument[categories-]{categories}
\externaldocument[topology-]{topology}
\externaldocument[sheaves-]{sheaves}
\externaldocument[sites-]{sites}
\externaldocument[stacks-]{stacks}
\externaldocument[fields-]{fields}
\externaldocument[algebra-]{algebra}
\externaldocument[brauer-]{brauer}
\externaldocument[homology-]{homology}
\externaldocument[derived-]{derived}
\externaldocument[simplicial-]{simplicial}
\externaldocument[more-algebra-]{more-algebra}
\externaldocument[smoothing-]{smoothing}
\externaldocument[modules-]{modules}
\externaldocument[sites-modules-]{sites-modules}
\externaldocument[injectives-]{injectives}
\externaldocument[cohomology-]{cohomology}
\externaldocument[sites-cohomology-]{sites-cohomology}
\externaldocument[dga-]{dga}
\externaldocument[dpa-]{dpa}
\externaldocument[sdga-]{sdga}
\externaldocument[hypercovering-]{hypercovering}
\externaldocument[schemes-]{schemes}
\externaldocument[constructions-]{constructions}
\externaldocument[properties-]{properties}
\externaldocument[morphisms-]{morphisms}
\externaldocument[coherent-]{coherent}
\externaldocument[divisors-]{divisors}
\externaldocument[limits-]{limits}
\externaldocument[varieties-]{varieties}
\externaldocument[topologies-]{topologies}
\externaldocument[descent-]{descent}
\externaldocument[perfect-]{perfect}
\externaldocument[more-morphisms-]{more-morphisms}
\externaldocument[flat-]{flat}
\externaldocument[groupoids-]{groupoids}
\externaldocument[more-groupoids-]{more-groupoids}
\externaldocument[etale-]{etale}
\externaldocument[chow-]{chow}
\externaldocument[intersection-]{intersection}
\externaldocument[pic-]{pic}
\externaldocument[weil-]{weil}
\externaldocument[adequate-]{adequate}
\externaldocument[dualizing-]{dualizing}
\externaldocument[duality-]{duality}
\externaldocument[discriminant-]{discriminant}
\externaldocument[derham-]{derham}
\externaldocument[local-cohomology-]{local-cohomology}
\externaldocument[algebraization-]{algebraization}
\externaldocument[curves-]{curves}
\externaldocument[resolve-]{resolve}
\externaldocument[models-]{models}
\externaldocument[functors-]{functors}
\externaldocument[equiv-]{equiv}
\externaldocument[pione-]{pione}
\externaldocument[etale-cohomology-]{etale-cohomology}
\externaldocument[proetale-]{proetale}
\externaldocument[relative-cycles-]{relative-cycles}
\externaldocument[more-etale-]{more-etale}
\externaldocument[trace-]{trace}
\externaldocument[crystalline-]{crystalline}
\externaldocument[spaces-]{spaces}
\externaldocument[spaces-properties-]{spaces-properties}
\externaldocument[spaces-morphisms-]{spaces-morphisms}
\externaldocument[decent-spaces-]{decent-spaces}
\externaldocument[spaces-cohomology-]{spaces-cohomology}
\externaldocument[spaces-limits-]{spaces-limits}
\externaldocument[spaces-divisors-]{spaces-divisors}
\externaldocument[spaces-over-fields-]{spaces-over-fields}
\externaldocument[spaces-topologies-]{spaces-topologies}
\externaldocument[spaces-descent-]{spaces-descent}
\externaldocument[spaces-perfect-]{spaces-perfect}
\externaldocument[spaces-more-morphisms-]{spaces-more-morphisms}
\externaldocument[spaces-flat-]{spaces-flat}
\externaldocument[spaces-groupoids-]{spaces-groupoids}
\externaldocument[spaces-more-groupoids-]{spaces-more-groupoids}
\externaldocument[bootstrap-]{bootstrap}
\externaldocument[spaces-pushouts-]{spaces-pushouts}
\externaldocument[spaces-chow-]{spaces-chow}
\externaldocument[groupoids-quotients-]{groupoids-quotients}
\externaldocument[spaces-more-cohomology-]{spaces-more-cohomology}
\externaldocument[spaces-simplicial-]{spaces-simplicial}
\externaldocument[spaces-duality-]{spaces-duality}
\externaldocument[formal-spaces-]{formal-spaces}
\externaldocument[restricted-]{restricted}
\externaldocument[spaces-resolve-]{spaces-resolve}
\externaldocument[formal-defos-]{formal-defos}
\externaldocument[defos-]{defos}
\externaldocument[cotangent-]{cotangent}
\externaldocument[examples-defos-]{examples-defos}
\externaldocument[algebraic-]{algebraic}
\externaldocument[examples-stacks-]{examples-stacks}
\externaldocument[stacks-sheaves-]{stacks-sheaves}
\externaldocument[criteria-]{criteria}
\externaldocument[artin-]{artin}
\externaldocument[quot-]{quot}
\externaldocument[stacks-properties-]{stacks-properties}
\externaldocument[stacks-morphisms-]{stacks-morphisms}
\externaldocument[stacks-limits-]{stacks-limits}
\externaldocument[stacks-cohomology-]{stacks-cohomology}
\externaldocument[stacks-perfect-]{stacks-perfect}
\externaldocument[stacks-introduction-]{stacks-introduction}
\externaldocument[stacks-more-morphisms-]{stacks-more-morphisms}
\externaldocument[stacks-geometry-]{stacks-geometry}
\externaldocument[moduli-]{moduli}
\externaldocument[moduli-curves-]{moduli-curves}
\externaldocument[examples-]{examples}
\externaldocument[exercises-]{exercises}
\externaldocument[guide-]{guide}
\externaldocument[desirables-]{desirables}
\externaldocument[coding-]{coding}
\externaldocument[obsolete-]{obsolete}
\externaldocument[fdl-]{fdl}
\externaldocument[index-]{index}

% Theorem environments.
%
\theoremstyle{plain}
\newtheorem{theorem}[subsection]{Theorem}
\newtheorem{proposition}[subsection]{Proposition}
\newtheorem{lemma}[subsection]{Lemma}

\theoremstyle{definition}
\newtheorem{definition}[subsection]{Definition}
\newtheorem{example}[subsection]{Example}
\newtheorem{exercise}[subsection]{Exercise}
\newtheorem{situation}[subsection]{Situation}

\theoremstyle{remark}
\newtheorem{remark}[subsection]{Remark}
\newtheorem{remarks}[subsection]{Remarks}

\numberwithin{equation}{subsection}

% Macros
%
\def\lim{\mathop{\mathrm{lim}}\nolimits}
\def\colim{\mathop{\mathrm{colim}}\nolimits}
\def\Spec{\mathop{\mathrm{Spec}}}
\def\Hom{\mathop{\mathrm{Hom}}\nolimits}
\def\Ext{\mathop{\mathrm{Ext}}\nolimits}
\def\SheafHom{\mathop{\mathcal{H}\!\mathit{om}}\nolimits}
\def\SheafExt{\mathop{\mathcal{E}\!\mathit{xt}}\nolimits}
\def\Sch{\mathit{Sch}}
\def\Mor{\mathop{\mathrm{Mor}}\nolimits}
\def\Ob{\mathop{\mathrm{Ob}}\nolimits}
\def\Sh{\mathop{\mathit{Sh}}\nolimits}
\def\NL{\mathop{N\!L}\nolimits}
\def\CH{\mathop{\mathrm{CH}}\nolimits}
\def\proetale{{pro\text{-}\acute{e}tale}}
\def\etale{{\acute{e}tale}}
\def\QCoh{\mathit{QCoh}}
\def\Ker{\mathop{\mathrm{Ker}}}
\def\Im{\mathop{\mathrm{Im}}}
\def\Coker{\mathop{\mathrm{Coker}}}
\def\Coim{\mathop{\mathrm{Coim}}}

% Boxtimes
%
\DeclareMathSymbol{\boxtimes}{\mathbin}{AMSa}{"02}

%
% Macros for moduli stacks/spaces
%
\def\QCohstack{\mathcal{QC}\!\mathit{oh}}
\def\Cohstack{\mathcal{C}\!\mathit{oh}}
\def\Spacesstack{\mathcal{S}\!\mathit{paces}}
\def\Quotfunctor{\mathrm{Quot}}
\def\Hilbfunctor{\mathrm{Hilb}}
\def\Curvesstack{\mathcal{C}\!\mathit{urves}}
\def\Polarizedstack{\mathcal{P}\!\mathit{olarized}}
\def\Complexesstack{\mathcal{C}\!\mathit{omplexes}}
% \Pic is the operator that assigns to X its picard group, usage \Pic(X)
% \Picardstack_{X/B} denotes the Picard stack of X over B
% \Picardfunctor_{X/B} denotes the Picard functor of X over B
\def\Pic{\mathop{\mathrm{Pic}}\nolimits}
\def\Picardstack{\mathcal{P}\!\mathit{ic}}
\def\Picardfunctor{\mathrm{Pic}}
\def\Deformationcategory{\mathcal{D}\!\mathit{ef}}
\usepackage{cleveref}
\crefname{originallemma}{original lemma}{original lemmas}
\Crefname{originallemma}{Original lemma}{Original lemmas}
\makeatletter
\newlabel{cohomology-lemma-perfect}{{08CQ}{source}{}{sourcepointer-1}{}}
\newlabel{cohomology-lemma-perfect@cref}{{[originallemma][0][]08CQ}{[1][0][]source}{}{}{}}
\makeatother
\crefname{originallemma}{original lemma}{original lemmas}
\Crefname{originallemma}{Original lemma}{Original lemmas}
\makeatletter
\newlabel{cohomology-lemma-finite-cohomology}{{08DN}{source}{}{sourcepointer-2}{}}
\newlabel{cohomology-lemma-finite-cohomology@cref}{{[originallemma][0][]08DN}{[1][0][]source}{}{}{}}
\makeatother
\crefname{originallemma}{original lemma}{original lemmas}
\Crefname{originallemma}{Original lemma}{Original lemmas}
\makeatletter
\newlabel{properties-lemma-finite-locally-free}{{05P2}{source}{}{sourcepointer-3}{}}
\newlabel{properties-lemma-finite-locally-free@cref}{{[originallemma][0][]05P2}{[1][0][]source}{}{}{}}
\makeatother
\crefname{originallemma}{original lemma}{original lemmas}
\Crefname{originallemma}{Original lemma}{Original lemmas}
\makeatletter
\newlabel{properties-lemma-quasi-coherent-colimit-finite-type}{{01PG}{source}{}{sourcepointer-4}{}}
\newlabel{properties-lemma-quasi-coherent-colimit-finite-type@cref}{{[originallemma][0][]01PG}{[1][0][]source}{}{}{}}
\makeatother
\crefname{originaldefinition}{original definition}{original definitions}
\Crefname{originaldefinition}{Original definition}{Original definitions}
\makeatletter
\newlabel{derived-definition-cone}{{014E}{source}{}{sourcepointer-5}{}}
\newlabel{derived-definition-cone@cref}{{[originaldefinition][0][]014E}{[1][0][]source}{}{}{}}
\makeatother
\crefname{originallemma}{original lemma}{original lemmas}
\Crefname{originallemma}{Original lemma}{Original lemmas}
\makeatletter
\newlabel{cohomology-lemma-cone-tor-amplitude}{{08CJ}{source}{}{sourcepointer-6}{}}
\newlabel{cohomology-lemma-cone-tor-amplitude@cref}{{[originallemma][0][]08CJ}{[1][0][]source}{}{}{}}
\makeatother
\title{A perfect object on a scheme with the resolution property is represented by a bounded complex of finite locally free modules}
\author{The Stacks Project Authors}
\date{stacks-project:perfect.tex:local-snapshot-20261012; GNU FDL 1.2 or later}
\begin{document}
\maketitle
\setcounter{subsection}{1}
\noindent\textit{Editorial note.} All delivered mathematical text is the author's own, copied from the pinned Stacks Project chapter file perfect.tex; no statement or proof step was rewritten, reordered or removed. The counted span is the author environment and its own proof environment. References inside the span are handed to the assembler as navigation stubs that carry the canonical tag of the cited tag only; whether each of them is load-bearing is recorded on the stub itself.\par\medskip
\noindent\textit{Allowed standard prerequisites (editorial).}\begin{itemize}
\item A perfect complex on a quasi-compact scheme has bounded tor amplitude and its cohomology modules are of finite type. invoked at proof lines 9800-9808 and again at proof lines 9820-9825. Provides the interval on which the induction runs and the finiteness used to build the surjection in the induction step.
\item A finitely presented module of tor dimension zero is finite locally free. invoked at proof lines 9812-9818. Identifies the top cohomology in the base case with a finite locally free module.
\item On a scheme with the resolution property the surjection onto a finite type module can be realized by a finite locally free module, and the kernel of such a surjection between finite locally free modules sits in an exact sequence controlled by the cone. invoked at proof lines 9825-9887. Produces the finite locally free module and the distinguished triangle used to reduce the interval.
\end{itemize}\medskip
\ifdefined\hypertarget\hypertarget{sourcepointer-1}{}\fi
\noindent\textit{Source reference.} Original lemma 08CQ: sources/stacks-chow/tags.txt, source lines 1--1. Cited at proof lines 9800-9802 to obtain the integers a and b bounding the tor amplitude; the induction set up on that interval is inside the counted span.\par\medskip
\ifdefined\hypertarget\hypertarget{sourcepointer-2}{}\fi
\noindent\textit{Source reference.} Original lemma 08DN: sources/stacks-chow/tags.txt, source lines 1--1. Cited twice, at proof lines 9812-9816 for the base case and at proof lines 9820-9825 for the finiteness of the top cohomology in the induction step; both uses are inside the counted span.\par\medskip
\ifdefined\hypertarget\hypertarget{sourcepointer-3}{}\fi
\noindent\textit{Source reference.} Original lemma 05P2: sources/stacks-chow/tags.txt, source lines 1--1. Invoked at proof lines 9816-9818 to conclude that the top cohomology in the base case is finite locally free; the base case argument is inside the counted span.\par\medskip
\ifdefined\hypertarget\hypertarget{sourcepointer-4}{}\fi
\noindent\textit{Source reference.} Original lemma 01PG: sources/stacks-chow/tags.txt, source lines 1--1. Used at proof lines 9826-9840 to realize the surjection onto the top cohomology by a finite locally free module, which is where the resolution property enters; the triangle built from it is inside the counted span.\par\medskip
\ifdefined\hypertarget\hypertarget{sourcepointer-5}{}\fi
\noindent\textit{Source reference.} Original definition 014E: sources/stacks-chow/tags.txt, source lines 1--1. Cited at proof lines 9851-9855 for the distinguished triangle whose cone has a shorter tor amplitude interval; the induction that uses it is inside the counted span.\par\medskip
\ifdefined\hypertarget\hypertarget{sourcepointer-6}{}\fi
\noindent\textit{Source reference.} Original lemma 08CJ: sources/stacks-chow/tags.txt, source lines 1--1. Invoked at proof lines 9856-9860 to see that the cone again has bounded tor amplitude, with a shorter interval; the application of the induction hypothesis to it is inside the counted span.\par\medskip
\begin{lemma}
\label{lemma-construct-strictly-perfect}
Let $X$ be a quasi-compact and quasi-separated scheme with the
resolution property.
Let $\mathcal{F}^\bullet$ be a bounded below complex of quasi-coherent
$\mathcal{O}_X$-modules representing a perfect object of
$D(\mathcal{O}_X)$. Then there exists a bounded complex
$\mathcal{E}^\bullet$ of finite locally free $\mathcal{O}_X$-modules
and a quasi-isomorphism $\mathcal{E}^\bullet \to \mathcal{F}^\bullet$.
\end{lemma}

\begin{proof}
Let $a, b \in \mathbf{Z}$ be integers such that $\mathcal{F}^\bullet$
has tor amplitude in $[a, b]$ and such that $\mathcal{F}^n = 0$ for
$n < a$. The existence of such a pair of integers
follows from Cohomology, Lemma \ref{cohomology-lemma-perfect}
and the fact that $X$ is quasi-compact.
If $b < a$, then $\mathcal{F}^\bullet$ is zero in the
derived category and the lemma holds.
We will prove by induction on
$b - a \geq 0$ that there exists a complex
$\mathcal{E}^a \to \ldots \to \mathcal{E}^b$
with $\mathcal{E}^i$ finite locally free
and a quasi-isomorphism $\mathcal{E}^\bullet \to \mathcal{F}^\bullet$.

\medskip\noindent
The base case is the case $b - a = 0$. In this case
$H^b(\mathcal{F}^\bullet) = H^a(\mathcal{F}^\bullet) =
\Ker(\mathcal{F}^a \to \mathcal{F}^{a + 1})$
is finite locally free. Namely, it is a finitely presented
$\mathcal{O}_X$-module of tor dimension $0$ and hence finite
locally free. See Cohomology, Lemmas \ref{cohomology-lemma-perfect} and
\ref{cohomology-lemma-finite-cohomology} and
Properties, Lemma \ref{properties-lemma-finite-locally-free}.
Thus we can take $\mathcal{E}^\bullet$ to be
$H^b(\mathcal{F}^\bullet)$ sitting in degree $b$.
The rest of the proof is dedicated to the induction step.

\medskip\noindent
Assume $b > a$. Observe that
$$
H^b(\mathcal{F}^\bullet) =
\Ker(\mathcal{F}^b \to \mathcal{F}^{b + 1})/
\Im(\mathcal{F}^{b - 1} \to \mathcal{F}^b)
$$
is a finite type quasi-coherent $\mathcal{O}_X$-module, see
Cohomology, Lemmas \ref{cohomology-lemma-perfect} and
\ref{cohomology-lemma-finite-cohomology}. Then we can find a finite type
quasi-coherent $\mathcal{O}_X$-module $\mathcal{F}$ and a map
$$
\mathcal{F} \longrightarrow
\Ker(\mathcal{F}^b \to \mathcal{F}^{b + 1})
$$
such that the composition with the projection onto
$H^b(\mathcal{F}^\bullet)$ is surjective.
Namely, we can write $\Ker(\mathcal{F}^b \to \mathcal{F}^{b + 1})$
as the filtered union of its finite type quasi-coherent submodules by
Properties, Lemma \ref{properties-lemma-quasi-coherent-colimit-finite-type}
and then one of these will do the job.
Next, we choose a finite locally free $\mathcal{O}_X$-module
$\mathcal{E}^b$ and a surjection $\mathcal{E}^b \to \mathcal{F}$ using
the resolution property of $X$. Consider the map of complexes
$$
\alpha : \mathcal{E}^b[-b] \to \mathcal{F}^\bullet
$$
and its cone $C(\alpha)^\bullet$, see
Derived Categories, Definition \ref{derived-definition-cone}.
We observe that $C(\alpha)^\bullet$ is nonzero only in degrees
$\geq a$, has tor amplitude in $[a, b]$ by
Cohomology, Lemma \ref{cohomology-lemma-cone-tor-amplitude},
and has $H^b(C(\alpha)^\bullet) = 0$ by construction.
Thus we actually find that $C(\alpha)^\bullet$ has tor amplitude
in $[a, b - 1]$. Hence the induction hypothesis applies to
$C(\alpha)^\bullet$ and we find a map of complexes
$$
(\mathcal{E}^a \to \ldots \to \mathcal{E}^{b - 1})
\longrightarrow
C(\alpha)^\bullet
$$
with properties as stated in the induction hypothesis. Unwinding
the definition of the cone this gives a commutative diagram
$$
\xymatrix{
\ldots \ar[r] &
\mathcal{E}^{b - 2} \ar[r] \ar[d] &
\mathcal{E}^{b - 1} \ar[r] \ar[d] &
0 \ar[r] \ar[d] &
\ldots \\
\ldots \ar[r] &
\mathcal{F}^{b - 2} \ar[r] &
\mathcal{F}^{b - 1} \oplus \mathcal{E}^b \ar[r] &
\mathcal{F}^b \ar[r] &
\ldots
}
$$
It is clear that we obtain a map of complexes
$(\mathcal{E}^a \to \ldots \to \mathcal{E}^b) \to \mathcal{F}^\bullet$.
We omit the verification that this map is a quasi-isomorphism.
\end{proof}

\end{document}
